Impact Chain Theory provides a structured framework for analysing climate-related risks by explicitly representing the relationships between hazard, exposure, vulnerability, and resulting impacts. Within this framework, risk is understood as the outcome of interacting system components rather than as a direct consequence of hazard conditions alone.

The European Drought Risk Atlas (EDORA) operationalises this approach through the construction of sector-specific impact chains, developed through a combination of literature review and expert consultation \cite{EDORA}. These impact chains describe sequences of processes linking hydrometeorological conditions to sectoral consequences. For instance, elevated temperatures and precipitation deficits may lead to soil moisture depletion, reduced crop productivity, and increased irrigation demand. Individual elements within these chains are categorised according to hazard, exposure, vulnerability, and impact.

To enable quantitative analysis, EDORA represents these conceptual relationships using proxy variables. Hazard indicators are derived from hydro-meteorological data, while drought impact indicators are obtained from datasets such as EDID (European Drought Impact Database). Vulnerability and exposure are approximated using socio-economic and environmental variables, including land use, economic activity, and climatic characteristics. These components are combined to estimate sectoral drought risk \cite{EDORA}. For example, EDORA uses impact chain theory to estimate the sectoral risk of wheat yield reduction. 

The resulting framework provides a formal representation in which drought impacts arise from the interaction between hazard conditions, exposed systems, and their underlying vulnerabilities. This relationship can be expressed in general functional form as:

\[
\text{I} = f(H, E, V)
\]

where $H$ denotes hazard conditions, $E$ represents exposure, and $V$ captures vulnerability. The function $f()$ describes how these components jointly determine the observed impacts, $I$.

In this thesis, the impact chain framework serves as the conceptual foundation for defining the machine learning input space. The relationship between hazard indicators and observed impacts is approximated using data-driven modelling techniques. This formulation allows the incorporation of heterogeneous data sources, including impact information derived from unstructured textual data. Examples of data that fit within this impact chain framework include meteorological, hydrological, and agricultural indices for hazard indicators. For exposure indicators, land-use ratios are primarily considered. For vulnerability, population density within a region (PopR) is used as an indicator.
